\documentclass[12pt,a4paper]{article}
\usepackage[english]{babel}
\usepackage{titlesec}

% \titleformat{\chapter} % command [display] % shape {\bfseries\huge} % format
%   {Chapter \thechapter} % label {0pt} % sep {} % before-code


\usepackage[T5]{fontenc}
\usepackage{mathptmx}[ptm]
\usepackage{a4wide, amssymb, epsfig, latexsym, array, hhline, fancyhdr}
\usepackage[normalem]{ulem}
%\usepackage{soul}
\usepackage{float}
\usepackage[makeroom]{cancel}
\usepackage{amsmath}
\usepackage{amsthm}
\usepackage{multicol, longtable, amscd}
\usepackage{diagbox}%Make diagonal lines in tables
\usepackage{booktabs}
\usepackage{alltt}
\usepackage[framemethod=tikz]{mdframed}% For highlighting paragraph backgrounds
\usepackage{indentfirst}
\usepackage{caption,subcaption}
\usepackage{svg}
\graphicspath{ {figures/} }
\usepackage{lastpage}
\usepackage[lined,boxed,commentsnumbered]{algorithm2e}
\usepackage{enumerate}
\usepackage{color}
\usepackage{graphicx}% Standard graphics package
\graphicspath{ {figures/} }
\usepackage{array}
\usepackage{tabularx, caption}
\usepackage{tabularray}
\usepackage{multirow}
\usepackage{multicol}
\usepackage{rotating}
\usepackage{graphics}
\usepackage{geometry}
\usepackage{setspace}
\usepackage{epsfig}
\usepackage{tikz}
\usepackage[numbers]{natbib}
\usepackage{subcaption}
\usepackage{dirtree}

\usetikzlibrary{arrows, snakes, backgrounds, calc}
\usepackage[unicode]{hyperref}
\hypersetup{urlcolor=blue,linkcolor=black,citecolor=black,colorlinks=true} 

\usepackage{xcolor}

\definecolor{codegreen}{rgb}{0,0.6,0}
\definecolor{codegray}{rgb}{0.5,0.5,0.5}
\definecolor{codepurple}{rgb}{0.58,0,0.82}
\definecolor{backcolour}{rgb}{0.95,0.95,0.92}

\usepackage{listings}
\lstdefinestyle{mystyle}{ backgroundcolor=\color{backcolour},   
    commentstyle=\color{codegreen}, keywordstyle=\color{magenta},
    numberstyle=\tiny\color{codegray}, stringstyle=\color{codepurple},
    basicstyle=\ttfamily\footnotesize, breakatwhitespace=false,         
    breaklines=true,                 
    captionpos=b,                    
    keepspaces=true,                 
    numbers=left,                    
    numbersep=5pt,                  
    showspaces=false,                
    showstringspaces=false, showtabs=false,                  
    tabsize=2 }
\lstset{style=mystyle}



\def\thesislayout{	% A4: 210 × 297
	\geometry{ a4paper, total={160mm,240mm},  % fix over page
		left=30mm, top=30mm, } }
\thesislayout

%\usepackage{fancyhdr}
\setlength{\headheight}{40pt}
\pagestyle{fancy}
\fancyhead{} % clear all header fields
\fancyhead[L]{
 \begin{tabular}{rl}
    \begin{picture}(25,15)(0,0) \put(0,-8){\includegraphics[width=10mm,
    height=10mm]{Images/hcmut.png}}
    %\put(0,-8){\epsfig{width=10mm,figure=hcmut.eps}}
   \end{picture}&
	%\includegraphics[width=8mm, height=8mm]{hcmut.png} & %
	\begin{tabular}{l}
        \textbf{ Ho Chi Minh City University of Technology } \\
        \textbf{ Faculty of Computer Science and Engineering }
	\end{tabular} 	
 \end{tabular}
} \fancyhead[R]{
	\begin{tabular}{l}
		\tiny \bf \\
		\tiny \bf \end{tabular}  } \fancyfoot{} % clear all footer fields
\fancyfoot[L]{\scriptsize Modern Speech Processing (CO5257) - Academic year 2024 - 2025}
\fancyfoot[R]{\scriptsize  Page {\thepage}/\pageref{LastPage}}
\renewcommand{\headrulewidth}{0.3pt}
\renewcommand{\footrulewidth}{0.3pt}

%%%%%%%%%%%%%%%%%%%%%%%%%%%%%%%%%%%%%%%%%%
% Tạo một môi trường mới cho các bảng đặc tả UC, giúp định dạng kiểu
\newenvironment{usecase_table}[1][Test]
{
\begin{table}[H]
% Left-Right Table padding
    \setlength{\tabcolsep}{6pt}
    \renewcommand{\arraystretch}{1.5}
    % \begin{tabularx}{\textwidth}{bs}
    \def\savedcaption{\caption{#1}}%
    \begin{tabular}{|>{\bfseries} p{0.2\linewidth}|p{0.7\linewidth}|}

} {
    \end{tabular}
    \savedcaption
\end{table}
}

\newenvironment{usecase_enum}
{
    \begin{enumerate*}[itemjoin={\newline}]
} {
    \end{enumerate*}
}

\makeatletter
\newenvironment{acknowledgments}{
	\small
	\begin{center}
	  {\bfseries Acknowledgments\vspace{-.5em}\vspace{\z@}}
	\end{center}
	\quotation
  }{}
\makeatother

% Declaration
\makeatletter
\newenvironment{declaration}{
	\small
	\begin{center}
	  {\bfseries Declaration\vspace{-.5em}\vspace{\z@}}
	\end{center}
	\quotation
  }{}
\makeatother

\makeatletter
\newenvironment{abstr}{
        \small
        \begin{center}
	{\bfseries Abstract \vspace{-.5em}\vspace{\z@}}
        \end{center}
        \quotation
    }{}
\makeatother

\makeatletter
\newenvironment{abstrsss}{
        \small
        \begin{center}
	{\bfseries Abstractss \vspace{-.5em}\vspace{\z@}}
        \end{center}
        \quotation
    }{}
\makeatother
%%%%%%%%%%%%%%%%%%%%%%%%%%%%%%%%%%%%%%%%%%%%%%5
\setcounter{secnumdepth}{4}
\setcounter{tocdepth}{3}
\makeatletter
\newcounter {subsubsubsection}[subsubsection]
\renewcommand\thesubsubsubsection{\thesubsubsection .\@alph\c@subsubsubsection}
\newcommand\subsubsubsection{\@startsection{subsubsubsection}{4}{\z@}%
                                     {-3.25ex\@plus -1ex \@minus -.2ex}%
                                     {1.5ex \@plus .2ex}%
                                     {\normalfont\normalsize\bfseries}}
\newcommand*\l@subsubsubsection{\@dottedtocline{3}{10.0em}{4.1em}}
\newcommand*{\subsubsubsectionmark}[1]{}
\makeatother



\everymath{\color{black}}%make in-line maths symbols blue to read/check easily

\sloppy
\captionsetup[figure]{labelfont={small,bf},textfont={small,it},belowskip=-1pt,aboveskip=-9pt}
%space removed between caption, figure, and text
\captionsetup[table]{labelfont={small,bf},textfont={small,it},belowskip=-1pt,aboveskip=7pt}
\setlength{\floatsep}{5pt plus 2pt minus 2pt}
\setlength{\textfloatsep}{5pt plus 2pt minus 2pt}
\setlength{\intextsep}{10pt plus 2pt minus 2pt}


% Declare a variable named Proc using \newcommand

\newcommand{\Uni}{Ho Chi Minh University - Vietnam National Unversity, Ho Chi Minh City}

\renewcommand{\labelitemi}{\normalfont\bfseries\textendash}
\renewcommand{\labelitemii}{\textbullet}
\renewcommand*\baselinestretch{1.5}\selectfont


\thesislayout

\begin{document}

\begin{titlepage}
    \begin{tikzpicture}[remember picture, overlay]
        \draw[line width = 4pt] ($(current page.north west) + (1.0in,-0.5in)$)
        rectangle ($(current page.south east) + (-0.4in,0.5in)$);
        \draw[line width=1.5pt]
        ($ (current page.north west) + (1.05in,-0.55in) $) rectangle ($ (current
        page.south east) + (-0.45in,0.55in) $);
    \end{tikzpicture}

    \begin{center}
        \large \textbf{VIETNAM NATIONAL UNIVERSITY, HO CHI MINH CITY} \\
        \large \textbf{HO CHI MINH UNIVERSITY OF TECHNOLOGY} \\
        \large \textbf{FACULTY OF COMPUTER SCIENCE AND ENGINEERING}
    \end{center}

    \begin{figure}[h!]
        \begin{center}
            \includegraphics[width=5cm]{Images/hcmut.png}
        \end{center}
    \end{figure}

    \begin{center}
        \begin{tabular}{c}
        \multicolumn{1}{l}{\textbf{{\Large Modern Speech Processing
        (CO5257)}}}\\
        ~~\\
        \hline
        \\
        \multicolumn{1}{l}{\textbf{{\Large Assignment Report}}}\\
        \\
        \textbf{{\Large Enhancing Vocoder Designs }}\\
        \textbf{{\Large with Emotion Embeddings and Emotional Cues}}\\
        \\
        \hline
        \end{tabular}
        \end{center}

    \vspace{1cm}
    \begin{table}[H]
        \begin{tabular}{rrl}
        \hspace{5 cm} & \textbf{Mentor}: & Nguyễn Đức Dũng\\
        
        & \textbf{Student}: & Trần Hà Tuấn Kiệt -- 2011493 \\
        
        \end{tabular}
        \end{table}
    \vspace{1cm}

    \begin{center}
        {\large Ho Chi Minh City, \the\month/\the\year}
    \end{center}
\end{titlepage}

\newpage

%\thispagestyle{empty} \chapter*{Work Assignment} \begin{table}[H] \centering
% \renewcommand{\arraystretch}{1.5} \begin{tabular}{|c|c|c|l|} \hline
% \textbf{No.}       & \textbf{Full name}                 & \textbf{Student ID}
% & \textbf{Work Assignment}                     \\
%         \hline %%%%%Student 1%%%%%%%%%%% \multirow{4}{*}{1} &
%         \multirow{4}{*}{Trần Hà Tuấn Kiệt} & \multirow{4}{*}{2011493} & -
%         Exploratory Data Analysis                  \\
%                            &                                    &
%                            & - System proposal and architecture           \\
%                            &                                    &
%                            & - Research and Develop with new technologies \\
%                            &                                    &
%                            & - Implement system endpoints                 \\
%         \hline %%%%%Student 2%%%%%%%%%%% \multirow{4}{*}{2} &
%         \multirow{4}{*}{Nguyễn Đức Thuỵ}   & \multirow{4}{*}{2012158} & -
%         Exploratory Data Analysis                  \\
%                            &                                    &
%                            & - System proposal and architecture           \\
%                            &                                    &
%                            & - Research and Develop with new technologies \\
%                            &                                    &
%                            & - Implement data processing flows            \\
%         \hline \end{tabular} \end{table} \newpage

% Revert title name back to English
\renewcommand{\contentsname}{Table of Contents}
\renewcommand{\listfigurename}{List of Figures}
\renewcommand{\listtablename}{List of Tables}
\renewcommand{\bibname}{References}
\renewcommand{\figurename}{Figure}

\tableofcontents
\newpage

%\thispagestyle{empty} \renewcommand\thechapter{Chapter \arabic{chapter}}

\section*{Abstract}
This article explores the integration of emotion embeddings and emotional cues
into vocoder architectures to enhance expressiveness in speech synthesis. Modern
vocoders, while capable of producing high-quality audio, often fail to
effectively capture the nuanced prosody and variability required for emotionally
rich speech. By conditioning vocoders with emotion-specific embeddings and
leveraging emotional cues such as pitch, rhythm, and intensity, this study aims
to bridge the gap between technical accuracy and human-like expressiveness.
Moreover, the study proposes a universal emotion-aware conditioning framework
applicable to neural, statistical, and hybrid vocoder designs. Through a
combination of theoretical exploration and empirical validation, the methodology
demonstrates improvements in emotional fidelity and naturalness across
synthesized speech outputs.
\newpage
\section{Introduction}

% Provide background on vocoder technology and its evolution.
Neural vocoders represent a cutting-edge technology in speech synthesis,
specializing in transforming intermediate acoustic representations, such as
spectrograms, into lifelike audio waveforms \cite{oord2016wavenet,
jang2021univnet, siuzdak2022vocos}. Renowned for their capability to produce
high-fidelity, natural-sounding speech, these models have become a pivotal
component of contemporary speech synthesis frameworks
\cite{albadawy2021vocbench, zarazaga2022neural}. Despite their significant
advances, a notable gap remains in achieving nuanced emotional expressiveness,
which is crucial for applications like virtual assistants, dubbing, and
human-computer interaction.

One of the primary shortcomings of existing vocoders lies in their inability to
accurately model the subtle variations in pitch, rhythm, intensity, and spectral
quality that characterize emotional speech \cite{siuzdak2022vocos}. Emotions
such as happiness, sadness, and anger often involve intricate acoustic
modulations, yet current models struggle to replicate these nuances. This
limitation stems partly from the design priorities of most vocoders, which focus
on naturalness and intelligibility rather than expressiveness.

Another challenge is the lack of generalization across speakers. Emotional
expressiveness varies significantly between individuals due to differences in
vocal characteristics and delivery styles. While vocoders can achieve high
performance when tailored to a specific speaker, their performance tends to
degrade when applied to other voices. This limitation arises because emotional
cues are often entangled with speaker-specific features, making it difficult for
models to adapt without extensive retraining \cite{zarazaga2022neural}.

The scarcity of high-quality, emotion-annotated datasets further exacerbates the
issue. Emotional speech synthesis requires large datasets that encompass a wide
range of emotions, speakers, and languages. However, collecting such datasets is
resource-intensive, and the subjective nature of emotion labeling introduces
inconsistencies. As a result, many vocoders are trained on datasets dominated by
neutral or mildly expressive speech, leaving them ill-equipped to handle more
dramatic or subtle emotional variations.


% Discuss the gap in existing vocoders regarding emotional expressiveness.
% Introduce the concept of emotion embeddings and their relevance. State the
% objectives: integrating emotion embeddings, improving emotional fidelity, and
% ensuring generalizability across vocoder types.


\newpage
\section {Methodology}


\newpage
\section {Experiments}

\newpage

\section{Evaluations}

\newpage
\section{Conclusion}

\newpage
\nocite{*}
\bibliographystyle{ieeetr}

\bibliography{refs}

\end{document}